% ---
% title: Choosing a portfolio
% tags: finance, combinations, complement
% ---

\begin{problem}
	\begin{context}
		A \emph{portfolio} is a collection of investments. Holding several
		stocks instead of one (\emph{diversification}) reduces risk. Here the
		portfolio is equally weighted, so only the \emph{set} of stocks matters.
	\end{context}
	An index has 30 stocks, 8 of them technology stocks. An analyst picks 5
	stocks. (a) How many portfolios are possible? (b) How many contain at least
	one tech stock? (c) How many contain exactly two tech stocks?
	(d) \emph{Lottery.} The Hungarian \emph{ötöslottó} draws 5 numbers out of
	90. How many different tickets are there?
\end{problem}

\begin{slideproblem}
	An index has 30 stocks, 8 of them technology stocks. An analyst
	picks 5 stocks (an equally weighted portfolio; order does not matter).
	(a) How many portfolios are possible?
	(b) How many contain at least one tech stock?
	(c) How many contain exactly two tech stocks?
\end{slideproblem}

\begin{slides}
	\small
	\exstatement
	\pause
	\sol
	\begin{enumerate}
		\item[(a)] Unordered, no repetition: $\binom{30}{5} =
			      \frac{30\cdot29\cdot28\cdot27\cdot26}{5!} = 142\,506$.
		\item[(b)] ``At least one'' $\to$ complement ``none'': all 5 from the 22
		      non-tech stocks, $\binom{22}{5}=26\,334$. So
		      $142\,506-26\,334 = 116\,172$.
		\item[(c)] Choose the tech part and the non-tech part separately, then
		      multiply: $\binom82\binom{22}{3} = 28\cdot1540 = 43\,120$.
	\end{enumerate}
	\pause
	\textbf{Lottery:} the Hungarian \emph{ötöslottó} draws 5 numbers out of 90,
	order irrelevant: $\binom{90}{5} = 43\,949\,268$ tickets --- one ticket wins
	the jackpot in 1 of about 44 million cases.
\end{slides}

\begin{solution}
	{\emergencystretch=1.5em\idea Unordered selections without repetition: binomial coefficients. ``At
	least one'' $\to$ complement ``none''. ``Exactly two'' $\to$ choose the two
	groups separately.\par}

	\step{(a)}
	\[ \binom{30}{5}=\frac{30\cdot29\cdot28\cdot27\cdot26}{5\cdot4\cdot3\cdot2\cdot1}
		=\frac{17\,100\,720}{120}=142\,506. \]

	\step{(b)} The opposite of ``at least one tech stock'' is ``no tech stock'':
	all 5 stocks come from the $30-8=22$ non-tech stocks, $\binom{22}{5}=26\,334$.
	Hence $142\,506-26\,334=116\,172$.

	\step{(c)} Choose 2 of the 8 tech stocks and 3 of the 22 others:
	\[ \binom82\binom{22}{3}=28\cdot1540=43\,120. \]

	\step{(d)} 5 numbers out of 90, order irrelevant:
	$\binom{90}{5}=43\,949\,268$.

	\begin{answer}
		(a) $142\,506$ \quad (b) $116\,172$ \quad (c) $43\,120$ \quad
		(d) $43\,949\,268$
	\end{answer}
	\begin{remark}
		\textbf{Common mistake in (b):} $8\cdot\binom{29}{4}$ (``choose one tech
		stock, then any 4 others''). This counts a portfolio with two tech stocks
		twice, once for each of them. Use the complement.\\
		\textbf{Lottery:} a single ticket wins the jackpot in 1 of about 44 million
		equally likely draws. Only part of the ticket revenue is paid back as
		prizes, so on average a ticket loses money (negative expected gain) --- we
		will compute such numbers in later weeks.
	\end{remark}
\end{solution}
